%------------------------------------------------------------------------------%
\begin{exercises}

\begin{exercise} Nos itens que seguem, encontre o volume do sólido obtido
  pela rotação da região limitada pelas curvas dadas em torno das retas
  especificadas.
\begin{alphalist}
  \item $y=2-\frac{1}{2}x$, $y=0$, $x=1$, $x=2$, em torno do eixo $x$.
  \item $y=x^3$, $y=0$, $x=2$, $x=20$, em torno do eixo $x$.
  \item $y=\sqrt{x-1}$, $y=0$, $x=1$, $x=5$, em torno do eixo $x$.
  \item $y=\frac{1}{x}$, $y=0$, $x=1$, $x=2$, em torno do eixo $x$.
  \item $y=x^3$, $y=x$, $x \geq 0$,  em torno do eixo $x$.
  \item $y=\sqrt{25-x^2}$, $y=0$, $x=2$, $x=4$, em torno do eixo $x$.
  \item $y^2=x$, $x=2y$, em torno do eixo $y$.
  \item $y=\sqrt{x}$, $y=x^2$, em torno do eixo $y=1$.
  \item $x=2\sqrt{y}$, $x=0$, $y=9$ em torno do eixo $y$.
  \item $y=\sin(x)$, $y=\cos(x)$, $0\leq x \leq \sfrac{\pi}{4}$
        em torno do eixo $y=-1$.
  \item $y=x^3$, $x=0$, $x=1$ em torno do eixo $x=2$.
  \item $y=5-x^2$, $y=\sfrac{1}{4}x^2$, $x=1$ em torno do eixo $x$.
\end{alphalist}
\begin{answer}
\begin{mathlist}
  \item \frac{19\pi}{12}
  \item \frac{128\pi}{7}
  \item 8\pi
  \item \frac{\pi}{2}
  \item \frac{5\pi}{21}
  \item \frac{94\pi}{3}
  \item \frac{64\pi}{15}
  \item \frac{11\pi}{30}
  \item 162\pi
  \item \left(2\sqrt{2}-\frac{3}{2}\right)\pi
  \item \frac{3\pi}{5}
  \item \frac{176\pi}{6}
\end{mathlist}
\end{answer}
\end{exercise}
%------------------------------------------------------------------------------%

%------------------------------------------------------------------------------%
\begin{exercise}
Determine o volume do sólido gerado bela rotação da região delimitada por
$y=\sqrt{x}$ e pelas retas $y=2$ e $x=0$ em torno
\begin{alphalist}
    \item do eixo x.
    \item do eixo y.
    \item da reta $y=2$.
    \item da reta $x=4$.
\end{alphalist}
\begin{answer}
\begin{mathlist}[2]
  \item 8\pi
  \item \frac{32\pi}{5}
  \item \frac{8\pi}{3}
  \item \frac{224\pi}{5}
\end{mathlist}
\end{answer}
\end{exercise}
%------------------------------------------------------------------------------%

%------------------------------------------------------------------------------%
\begin{exercise}
Considere a região $\mathcal{R}$ do primeiro quadrante delimitada pelas
curvas $y=-x^2+4x$ e $y=x$.
\begin{alphalist}
  \item Esboce a região $\mathcal{R}$ e calcule a sua área.
  \item Encontre o volume do sólido gerado pela rotação da região $\mathcal{R}$ em torno do eixo $x$.
  \item Encontre o volume do sólido gerado pela rotação da região $\mathcal{R}$ em torno da reta $y=-1$.
  \item Encontre o volume do sólido gerado pela rotação da região $\mathcal{R}$ em torno da reta $y=3$.
\end{alphalist}
\begin{answer}
\begin{mathlist}[2]
  \item \frac{9}{2}
  \item \frac{108 \pi}{5}
  \item \frac{153 \pi}{5}
  \item \frac{27  \pi}{5}
\end{mathlist}
\end{answer}
\end{exercise}
%------------------------------------------------------------------------------%

%------------------------------------------------------------------------------%
\begin{exercise}
Considere a região $\mathcal{R}$ do primeiro quadrante delimitada pelas
curvas $y=4-x^2$ e $y=2-x$.
\begin{alphalist}
  \item Esboce a região $\mathcal{R}$ e calcule a sua área.
  \item Encontre o volume do sólido gerado pela rotação da região $\mathcal{R}$ em torno do eixo $x$.
  \item Encontre o volume do sólido gerado pela rotação da região $\mathcal{R}$ em torno da reta $y=-1$.
  \item Encontre o volume do sólido gerado pela rotação da região $\mathcal{R}$ em torno da reta $y=5$.
\end{alphalist}
\begin{answer}
\begin{mathlist}[2]
  \item \frac{9}{2}
  \item \frac{108 \pi}{5}
  \item \frac{153 \pi}{5}
  \item \frac{117 \pi}{5}
\end{mathlist}
\end{answer}
\end{exercise}
%------------------------------------------------------------------------------%

%------------------------------------------------------------------------------%
\begin{exercise}
Considere a região $\mathcal{R}$ do primeiro quadrante delimitada pelas curvas $y=x$ e $y=\sqrt{x}$.
\begin{alphalist}
  \item Esboce a região $\mathcal{R}$ e calcule a sua área.
  \item Encontre o volume do sólido gerado pela rotação da região $\mathcal{R}$ em torno do eixo $x$.
  \item Encontre o volume do sólido gerado pela rotação da região $\mathcal{R}$ em torno do eixo $y$.
  \item Encontre o volume do sólido gerado pela rotação da região $\mathcal{R}$ em torno da reta $y=-1$.
  \item Encontre o volume do sólido gerado pela rotação da região $\mathcal{R}$ em torno da reta $y=1$.
  \item Encontre o volume do sólido gerado pela rotação da região $\mathcal{R}$ em torno da reta $x=1$.
\end{alphalist}
\begin{answer}
\begin{mathlist}[3]
  \item \frac{9}{2}
  \item \frac{  \pi}{6}
  \item \frac{2 \pi}{15}
  \item \frac{  \pi}{2}
  \item \frac{  \pi}{6}
  \item \frac{  \pi}{5}
\end{mathlist}
\end{answer}
\end{exercise}
%------------------------------------------------------------------------------%

%------------------------------------------------------------------------------%
\begin{exercise}
Esse exercício tem como objetivo usar integrais para demonstrar a
fórmula de volume de alguns sólidos.
Em cada item, encontre o volume do solido descrito.
\begin{alphalist}
  \item uma pirâmide de base quadrada com lado $L$ e cuja altura é $h$.
  \item Um cone circular reto de altura $h$ e raio da base $r$.
  \item Um tronco de cone circular reto de altura $h$, raio da base inferior
        $R$ e  raio da base superior $r$.
  \item Uma calota de uma esfera de raio $r$ e altura $h$.
\end{alphalist}
\begin{answer}
\begin{mathlist}[2]
  \item \frac{L^2h}{3}
  \item \frac{1}{3}r^2h
  \item \pi h \left(R^2 + Rh+ r^2 \right)
  \item \pi h^2 \left(r -\frac{h}{3} \right)
\end{mathlist}
\end{answer}
\end{exercise}
%------------------------------------------------------------------------------%

%------------------------------------------------------------------------------%
%\begin{exercise}
%\begin{answer}
%\end{answer}
%\end{exercise}

%------------------------------------------------------------------------------%
%\begin{exercise}
%\begin{mathlist}[2]
%  \item
%  \item
%  \item
%\end{mathlist}
%\begin{answer}
%  \begin{alphalist}[3]
%    \item
%    \item
%    \item
%  \end{alphalist}
%\end{answer}
%\end{exercise}

\end{exercises}
%------------------------------------------------------------------------------%
